\documentclass[10pt,a4paper]{amsart}
\usepackage[margin=19mm]{geometry}
\usepackage{amsmath,amssymb,mathtools}
\usepackage[scr=rsfs,cal=euler]{mathalfa}
\usepackage{xfrac}
\usepackage[unicode,hidelinks,bookmarks=false]{hyperref}
\newcommand{\ie}{i.e.\ }
\newcommand{\cf}{cf.\ }
\newcommand{\resp}{respectively\ }
\newcommand{\Subsection}{\subsection}
\providecommand{\nobreakdash}{\nobreak\hbox{-}}
\setlength{\emergencystretch}{2em}
\multlinegap0pt
\usepackage{float}
\usepackage{amssymb}
\newtheorem{theorem}{Theorem}
\newtheorem{corollary}{Corollary}
\title{DISTORTION OF THE TRIANGULAR RATIO METRIC UNDER MOEBIUS TRANSFORMATIONS}
\author{S. Nasyrov}
\date{Source: arXiv:2605.25779v3 (CC BY 4.0)}
\begin{document}
\maketitle

\noindent\emph{Source and licence.} DISTORTION OF THE TRIANGULAR RATIO METRIC UNDER MOEBIUS TRANSFORMATIONS. Version arXiv:2605.25779v3 (CC BY 4.0), distributed by arXiv under the Creative Commons Attribution 4.0 International licence, http://creativecommons.org/licenses/by/4.0/, as recorded in the arXiv licence metadata for this exact version and shown on the abstract page https://arxiv.org/abs/2605.25779v3. Full licence notice: \texttt{licenses/arxiv-2605.25779-CC-BY-4.0.txt}.\par\smallskip

\noindent\textit{Editorial scope.} Counted unit: the article's corollary cor3 (source0.tex lines 296--305). The article's complete original proof of it is delivered with it (source0.tex lines 307--330, one \texttt{proof} environment of 24 lines). No mathematics is written, repaired or re-derived: the statement and the proof are the article's own lines, in the article's own notation, and no retained line is substituted or re-flowed.  Retained uncounted as the source-local context the proof needs: lines 72--74; lines 83--89; lines 85--87; lines 186--188; lines 199--201; lines 203--205; lines 221--223.  Nothing was omitted from any retained span, so the package carries no omission record.  No retained line contains a citation, so the wrapper carries no bibliography and no bibliography entry.  No retained line contains a figure environment, an image-inclusion command or drawing code, so no image is delivered and the package declares no asset dependency.  Editorial formatting only, in addition: an AMS \texttt{amsart} preamble that restates the article's own declarations and macros used by the retained lines, namely the environments \texttt{theorem}, \texttt{corollary} on separate counters, together with the standard packages those lines need.  The retained environments print in this document as Theorem 1, Corollary 1 in order of appearance, whereas the article prints the same items with the numbering of its own full document.  The complete native source of the article as distributed in the arXiv e-print for this exact version is bundled unchanged as \texttt{sources/201377/source0.tex}.
\begin{equation}\label{moeb}
w=f(z)=\frac{z+a}{1+za}, \quad 0\le a<1,
\end{equation}

\begin{theorem}\label{th1}
Let $w=f(z)$ be a Moebius automorphism of the unit disk $\mathbb{U}$  of the form \eqref{moeb}. Then for the triangular ratio metric  $s_\mathbb{U}$ in $\mathbb{U}$ and every points $z_1$, $z_2\in \mathbb{U}$ we have
\begin{equation}\label{main}
s_\mathbb{U}(w_1,w_2)\le (1+a)s_\mathbb{U}(z_1,z_2), \quad w_1=f(z_1), w_2=f(z_2),
\end{equation}
and the constant $1+a$ is the best possible.
\end{theorem}

\begin{equation}
s_\mathbb{U}(w_1,w_2)\le (1+a)s_\mathbb{U}(z_1,z_2), \quad w_1=f(z_1), w_2=f(z_2),
\end{equation}

\begin{equation}\label{2a}
s_\mathbb{U}(w_1,w_2)\le s_\mathbb{U}(z_1,z_2),
\end{equation}

\begin{equation}\label{R1}
|(1-R)e^{i\theta}+1/a|=R.
\end{equation}

\begin{equation}\label{R}
R\ge \frac{1}{2}\left(1+\frac{1}{a}\right).
\end{equation}

\begin{equation}\label{1a}
1+a\ge\frac{2R}{2R-1}.
\end{equation}

\begin{corollary}\label{cor3}
Let the assumptions of Theorem~\ref{th1}  hold and $e^{i\phi}$ be the contact point  of the unit circle and the maximal ellipse with foci at the points $w_1$ and $w_2$, contained in~$\mathbb{U}$. If $\cos\phi\ge a$, then
$$
s_\mathbb{U}(w_1,w_2)\le \left(1+a\,\frac{\cos\phi-a}{1-a\cos\phi}\right)\,s_\mathbb{U}(z_1,z_2),
$$
and if $\cos\phi< a$, then
$$
s_\mathbb{U}(w_1,w_2)\le s_\mathbb{U}(z_1,z_2).
$$
\end{corollary}

\begin{proof}
In the proof of Theorem~\ref{th1} we use estimate \eqref{R} and its corollary \eqref{1a} for the radius $R$. But if we know the contact point $e^{i\phi}$, we can  use the exact value of $R$. From \eqref{R1} we find
$$
R=\frac{1+2a\cos\theta+a^2}{2a(a+\cos \theta) }\,.
$$
where $e^{i\theta}=f^{-1}(e^{i\varphi})$.
Then, instead of $(1+a)$ in \eqref{main}, we obtain the constant
$$
\frac{2R}{2R-1}=1+a\,\frac{a+\cos\theta}{1+a\cos \theta}\,.
$$
Since
$$
e^{i\theta}=\frac{e^{i\phi}-a}{1-e^{i\phi }a}\,,
$$
we obtain
$$
\cos\theta= \frac{(1+a^2)\cos\phi-2a}{1+a^2-2a\cos \phi}\,,
$$
therefore,
$$
\frac{2R}{2R-1}=1+a\,\frac{\cos\phi-a}{1-a\cos\phi}\,.
$$
In the case $\cos\phi< a$ we have the external tangency of circles. In the proof above it was showed that then \eqref{2a} holds.
\end{proof}

\end{document}
